\section{\method{} Audit Algorithm}

\begin{algorithm}[t]
\caption{\method{} black-box deletion audit}
\begin{algorithmic}[1]
\STATE Build or import provenance graph $G$
\STATE Estimate target priorities $\pi(z)$ from influence and provenance degree
\STATE Generate direct, paraphrase, semantic, retrieval, bridge, and extraction probes
\STATE Allocate probe budget using $\pi(z)$ and a minimum per-target budget
\FOR{each selected probe $q$ for target $z$}
  \STATE Query post-deletion system $S_-$
  \STATE Score evidence vector $e(S_-,q,z)$
\ENDFOR
\STATE Compute per-target and aggregate risk UCBs
\STATE Evaluate retain utility on $Q_{\mathrm{ret}}$
\IF{any probe exceeds violation tolerance}
  \STATE \textbf{return} fail
\ELSIF{aggregate UCB is below tolerance and retain utility is acceptable}
  \STATE \textbf{return} pass
\ELSE
  \STATE \textbf{return} escalate
\ENDIF
\end{algorithmic}
\end{algorithm}

The audit is intentionally black-box: it can be applied when the operator has access to system responses and retrieval metadata but not necessarily to training gradients or model weights. This design matches many practical foundation-model deployments where an application team controls the RAG pipeline and adapters but depends on external base models.

The algorithm uses two efficiency levers. First, provenance removes irrelevant audit targets by tracing deletion records to downstream artifacts. Second, valuation-guided allocation spends more probes on targets with high pre-deletion influence or many derived artifacts. In the current prototype, influence is represented by target metadata and provenance degree; the full benchmark can replace this with retrieval frequency, answer attribution, or Shapley-style estimates.

\begin{figure*}[t]
\centering
\includegraphics[width=0.95\textwidth]{../figures/adcu_pdf/adcu_pipeline.pdf}
\caption{ADCU audit pipeline. Deletion requests are expanded through provenance, prioritized by influence and artifact degree, audited through black-box probes, and summarized as risk confidence bounds plus retain-utility checks.}
\label{fig:adcu_pipeline}
\end{figure*}
